\subsection{Lecture Students and Unfamiliar Visitors}\label{sec:audience-encoding}

Student status is not a movement class. The implementation separates initial exit knowledge, entrance preference, age, response and staffing. The brief assumes regular lecture students know emergency exits and have no Crowd Managers \cite{brief}; it does not establish a faster walking speed or a precise response distribution.

Let $k$ denote \texttt{knowledge}. The foyer is known initially; each other exit eligible for the person's seating level is included in their initial known set independently with probability $k$, using a fixed person--exit draw. This is route knowledge, not awareness of remote queues or closures. Let $p_H$ denote \texttt{familiarityShare}. That legacy field name is the share with the 12 s foyer utility preference described above, not a measure of knowing the building well. A well-informed regular attendee can still prefer their usual entrance; an unfamiliar visitor can know the entrance without knowing alternatives. The two controls are independent.

The legacy lecture preset uses $k=1$, $p_H=0$, no managers, 10\% older share and recognition and preparation medians of 3/8 s. The concert reference uses $k=1$, $p_H=0.35$, six managers, 20\% older share and 6/14 s medians. Both initially know all eligible exits. Their difference is a bundled scenario assumption, not a measured student effect or an isolated familiarity comparison. It is retained for reproducibility and labelled accordingly in the interface.

Three controlled information presets now keep demographics, speeds, response, groups and compliance fixed, with no managers: $(k,p_H)=(1,0.35)$, $(0.5,0.35)$, and $(1,0)$. These isolate reduced exit knowledge and removal of foyer bias respectively. The value 0.5 is a sensitivity assumption, not an estimated visitor probability. Tests verify that a common seed preserves the other traits. No new lecture-versus-visitor Monte Carlo result is claimed in this review.

Tong et al.'s observations in a university library show that familiarity and exit choice depend on context, including whether an evacuation is a drill \cite{tong2025}. They do not calibrate these Bailey probabilities. Another limitation remains: knowing an exit name implies graph-route knowledge in the current model. It does not represent incomplete cognitive maps, wandering, or interpretation of successive arrows. This is particularly relevant to unfamiliar visitors and must be resolved before treating the sign results as a real wayfinding prediction.
